% --- Archon physics formalization source begin ---
% archon:physics
% archon:covers PhyXMiniProblems/problem_phyx_mini_0693.lean
% archon:source-report reports/phyx_mini/problem_phyx_mini_0693.source.json

\chapter{Physics problem phyx\_mini\_0693}
\label{ch:PhyXMiniProblems_problem_phyx_mini_0693}

\paragraph{Problem source.}
\#\# Physical scenario

An outfielder throws a baseball to his catcher in an attempt to throw out a runner at home plate. The ball bounces once before reaching the catcher. Assume the angle at which the bounced ball leaves the ground is the same as the angle at which the outfielder threw it as shown in figure, but that the ball's speed after the bounce is one-half of what it was before the bounce.

\#\# Question

Determine the ratio of the time interval for the one-bounce throw to the flight time for the no-bounce throw.

\#\# Answer choices

- A: A: 0.870

- B: B: 0.647

- C: C: 0.949

- D: D: 0.511

\#\# Figure caption (auxiliary; use the image as primary evidence)

The image depicts a physics diagram illustrating projectile motion with three trajectories. 

1. **Trajectories:**
   - The paths are shown with dashed lines. The first path, in blue, represents a lower angle trajectory. The second path, also blue, shows a higher angle trajectory. The third path, in green, represents the highest and longest trajectory at an angle of 45 degrees.

2. **Angles:**
   - Two angles, \textbackslash{}( \textbackslash{}theta \textbackslash{}) and 45 degrees, are marked. \textbackslash{}( \textbackslash{}theta \textbackslash{}) represents the launch angles for the first two trajectories.

3. **Vectors:**
   - Red arrows depict the initial velocity vectors. One vector is directed at an angle \textbackslash{}( \textbackslash{}theta \textbackslash{}) for the first lower trajectory, and another red arrow at the same angle \textbackslash{}( \textbackslash{}theta \textbackslash{}) for the second trajectory starting from the peak of the first arc. For the trajectory at 45 degrees, the arrow indicates a higher angle of launch.

4. **Labels:**
   - The distance along the horizontal axis is labeled \textbackslash{}( D \textbackslash{}) with a double-headed arrow indicating the range of the projectile.

This diagram is likely illustrating the concept of projectile motion showing how different launch angles affect the range and height of the projectile.

\#\# Dataset metadata

Kinematics; Multi-Formula Reasoning, Spatial Relation Reasoning

\paragraph{Recorded answer/context.}
C: C: 0.949

\paragraph{Figure/image path.}
/root/proposal\_for\_physic/science-mango/phyx\_mini\_run/phyx\_data/test\_image/693.png

\paragraph{Formalization target.}
create a compiling Lean file with sorry bodies at `PhyXMiniProblems/problem\_phyx\_mini\_0693.lean`. The Lean declarations must preserve the physical quantities, dimensions or dimensional roles, figure labels, governing-law hypotheses, and final relation expressed by this problem.
Use Mathlib/Physlib names found through LeanExplore where available. If a physics API is missing, introduce faithful local abstractions rather than scalar placeholder aliases.

\begin{theorem}[Physics formalization target]
\label{thm:physics:phyx_mini_0693:target}
\lean{PhyXMiniProblems.ProblemPhyXMini0693.problem_phyx_mini_0693}
\uses{def:physics:phyx-mini-0693:phyxminiproblems-problemphyxmini0693-baseballonebouncesetup, def:physics:phyx-mini-0693:phyxminiproblems-problemphyxmini0693-matchesprimarybaseballfigure, def:physics:phyx-mini-0693:phyxminiproblems-problemphyxmini0693-matchesbaseballthrowscenario, def:physics:phyx-mini-0693:phyxminiproblems-problemphyxmini0693-hasphysicalbaseballthrowparameters, def:physics:phyx-mini-0693:phyxminiproblems-problemphyxmini0693-satisfiesideallevelgroundprojectilelaws, def:physics:phyx-mini-0693:phyxminiproblems-problemphyxmini0693-onebouncetonobouncetimeratio, def:physics:phyx-mini-0693:phyxminiproblems-problemphyxmini0693-matchesdisplayedthreedecimalratio, def:physics:phyx-mini-0693:phyxminiproblems-problemphyxmini0693-isuniqueclosestdisplayedratio}
The assigned autoformalize agent should translate this physics problem into Lean declarations in the covered file, with theorem and lemma proof bodies written as `by sorry`.
\end{theorem}
\begin{proof}
This is an autoformalization task, not a proof task. Produce statements that can later be proved by the physics prover without weakening the physical model.
\end{proof}

% --- Archon declaration topology begin ---
\section{Declaration topology}
\begin{definition}[acceleration Dimension]
  \label{def:physics:phyx-mini-0693:phyxminiproblems-problemphyxmini0693-accelerationdimension}
  \lean{PhyXMiniProblems.ProblemPhyXMini0693.accelerationDimension}
  The physical dimension of acceleration, L T⁻²
\end{definition}
\begin{proof}
  This is the declared model definition of acceleration Dimension.
\end{proof}
\begin{definition}[Length Quantity]
  \label{def:physics:phyx-mini-0693:phyxminiproblems-problemphyxmini0693-lengthquantity}
  \lean{PhyXMiniProblems.ProblemPhyXMini0693.LengthQuantity}
  A nonnegative physical length
\end{definition}
\begin{proof}
  This is the declared model definition of Length Quantity.
\end{proof}
\begin{definition}[Duration Quantity]
  \label{def:physics:phyx-mini-0693:phyxminiproblems-problemphyxmini0693-durationquantity}
  \lean{PhyXMiniProblems.ProblemPhyXMini0693.DurationQuantity}
  A nonnegative physical duration
\end{definition}
\begin{proof}
  This is the declared model definition of Duration Quantity.
\end{proof}
\begin{definition}[Acceleration Quantity]
  \label{def:physics:phyx-mini-0693:phyxminiproblems-problemphyxmini0693-accelerationquantity}
  \lean{PhyXMiniProblems.ProblemPhyXMini0693.AccelerationQuantity}
  \uses{def:physics:phyx-mini-0693:phyxminiproblems-problemphyxmini0693-accelerationdimension}
  A nonnegative magnitude of gravitational acceleration
\end{definition}
\begin{proof}
  This is the declared model definition of Acceleration Quantity.
\end{proof}
\begin{definition}[nonnegative SIReadout]
  \label{def:physics:phyx-mini-0693:phyxminiproblems-problemphyxmini0693-nonnegativesireadout}
  \lean{PhyXMiniProblems.ProblemPhyXMini0693.nonnegativeSIReadout}
  Read a nonnegative dimensionful quantity in coherent SI units
\end{definition}
\begin{proof}
  This is the declared model definition of nonnegative SIReadout.
\end{proof}
\begin{definition}[length In Meters]
  \label{def:physics:phyx-mini-0693:phyxminiproblems-problemphyxmini0693-lengthinmeters}
  \lean{PhyXMiniProblems.ProblemPhyXMini0693.lengthInMeters}
  \uses{def:physics:phyx-mini-0693:phyxminiproblems-problemphyxmini0693-lengthquantity, def:physics:phyx-mini-0693:phyxminiproblems-problemphyxmini0693-nonnegativesireadout}
  Metre readout of a physical length
\end{definition}
\begin{proof}
  This is the declared model definition of length In Meters.
\end{proof}
\begin{definition}[duration In Seconds]
  \label{def:physics:phyx-mini-0693:phyxminiproblems-problemphyxmini0693-durationinseconds}
  \lean{PhyXMiniProblems.ProblemPhyXMini0693.durationInSeconds}
  \uses{def:physics:phyx-mini-0693:phyxminiproblems-problemphyxmini0693-durationquantity, def:physics:phyx-mini-0693:phyxminiproblems-problemphyxmini0693-nonnegativesireadout}
  Second readout of a physical duration
\end{definition}
\begin{proof}
  This is the declared model definition of duration In Seconds.
\end{proof}
\begin{definition}[speed In Meters Per Second]
  \label{def:physics:phyx-mini-0693:phyxminiproblems-problemphyxmini0693-speedinmeterspersecond}
  \lean{PhyXMiniProblems.ProblemPhyXMini0693.speedInMetersPerSecond}
  Metre-per-second readout of a physical speed
\end{definition}
\begin{proof}
  This is the declared model definition of speed In Meters Per Second.
\end{proof}
\begin{definition}[acceleration In Meters Per Second Squared]
  \label{def:physics:phyx-mini-0693:phyxminiproblems-problemphyxmini0693-accelerationinmeterspersecondsquared}
  \lean{PhyXMiniProblems.ProblemPhyXMini0693.accelerationInMetersPerSecondSquared}
  \uses{def:physics:phyx-mini-0693:phyxminiproblems-problemphyxmini0693-accelerationquantity, def:physics:phyx-mini-0693:phyxminiproblems-problemphyxmini0693-nonnegativesireadout}
  Metre-per-second-squared readout of gravitational acceleration
\end{definition}
\begin{proof}
  This is the declared model definition of acceleration In Meters Per Second Squared.
\end{proof}
\begin{definition}[Throw Segment]
  \label{def:physics:phyx-mini-0693:phyxminiproblems-problemphyxmini0693-throwsegment}
  \lean{PhyXMiniProblems.ProblemPhyXMini0693.ThrowSegment}
  The complete direct flight and the two arcs of the one-bounce flight
\end{definition}
\begin{proof}
  This is the declared model definition of Throw Segment.
\end{proof}
\begin{definition}[Ground Point]
  \label{def:physics:phyx-mini-0693:phyxminiproblems-problemphyxmini0693-groundpoint}
  \lean{PhyXMiniProblems.ProblemPhyXMini0693.GroundPoint}
  Ground points distinguished by the diagram and prose
\end{definition}
\begin{proof}
  This is the declared model definition of Ground Point.
\end{proof}
\begin{definition}[Trajectory Color]
  \label{def:physics:phyx-mini-0693:phyxminiproblems-problemphyxmini0693-trajectorycolor}
  \lean{PhyXMiniProblems.ProblemPhyXMini0693.TrajectoryColor}
  Colors used to distinguish the trajectories in the supplied image
\end{definition}
\begin{proof}
  This is the declared model definition of Trajectory Color.
\end{proof}
\begin{definition}[Figure Angle Label]
  \label{def:physics:phyx-mini-0693:phyxminiproblems-problemphyxmini0693-figureanglelabel}
  \lean{PhyXMiniProblems.ProblemPhyXMini0693.FigureAngleLabel}
  The two angle labels printed in the image
\end{definition}
\begin{proof}
  This is the declared model definition of Figure Angle Label.
\end{proof}
\begin{definition}[Baseball Throw Figure]
  \label{def:physics:phyx-mini-0693:phyxminiproblems-problemphyxmini0693-baseballthrowfigure}
  \lean{PhyXMiniProblems.ProblemPhyXMini0693.BaseballThrowFigure}
  \uses{def:physics:phyx-mini-0693:phyxminiproblems-problemphyxmini0693-throwsegment, def:physics:phyx-mini-0693:phyxminiproblems-problemphyxmini0693-groundpoint, def:physics:phyx-mini-0693:phyxminiproblems-problemphyxmini0693-trajectorycolor, def:physics:phyx-mini-0693:phyxminiproblems-problemphyxmini0693-figureanglelabel}
  Schematic evidence read from the primary image. The coordinates are drawing coordinates rather than measured physical distances; the physical baseline length is stored separately as distanceD below
\end{definition}
\begin{proof}
  This is the declared model definition of Baseball Throw Figure.
\end{proof}
\begin{definition}[Baseball One Bounce Setup]
  \label{def:physics:phyx-mini-0693:phyxminiproblems-problemphyxmini0693-baseballonebouncesetup}
  \lean{PhyXMiniProblems.ProblemPhyXMini0693.BaseballOneBounceSetup}
  \uses{def:physics:phyx-mini-0693:phyxminiproblems-problemphyxmini0693-lengthquantity, def:physics:phyx-mini-0693:phyxminiproblems-problemphyxmini0693-durationquantity, def:physics:phyx-mini-0693:phyxminiproblems-problemphyxmini0693-accelerationquantity, def:physics:phyx-mini-0693:phyxminiproblems-problemphyxmini0693-throwsegment, def:physics:phyx-mini-0693:phyxminiproblems-problemphyxmini0693-baseballthrowfigure}
  Independent physical quantities for the direct throw and both one-bounce segments. In particular, no duration or time ratio is defined from an answer choice. Relations among these quantities are supplied by scenario and governing-law premises below
\end{definition}
\begin{proof}
  This is the declared model definition of Baseball One Bounce Setup.
\end{proof}
\begin{definition}[Matches Primary Baseball Figure]
  \label{def:physics:phyx-mini-0693:phyxminiproblems-problemphyxmini0693-matchesprimarybaseballfigure}
  \lean{PhyXMiniProblems.ProblemPhyXMini0693.MatchesPrimaryBaseballFigure}
  \uses{def:physics:phyx-mini-0693:phyxminiproblems-problemphyxmini0693-baseballonebouncesetup}
  Qualitative and spatial facts visible in the supplied raster image
\end{definition}
\begin{proof}
  This is the declared model definition of Matches Primary Baseball Figure.
\end{proof}
\begin{definition}[Matches Baseball Throw Scenario]
  \label{def:physics:phyx-mini-0693:phyxminiproblems-problemphyxmini0693-matchesbaseballthrowscenario}
  \lean{PhyXMiniProblems.ProblemPhyXMini0693.MatchesBaseballThrowScenario}
  \uses{def:physics:phyx-mini-0693:phyxminiproblems-problemphyxmini0693-lengthinmeters, def:physics:phyx-mini-0693:phyxminiproblems-problemphyxmini0693-speedinmeterspersecond, def:physics:phyx-mini-0693:phyxminiproblems-problemphyxmini0693-baseballonebouncesetup}
  Problem-statement data and the comparison convention implicit in the figure. The no-bounce throw and the first leg of the one-bounce throw start with the same speed. The second blue leg leaves the bounce with half the incoming impact speed and at the same angle θ as the first blue leg. The direct green throw is at 45°. The ideal-motion laws below separately relate the incoming impact speed to the first launch speed. The range equations state only that both routes connect the same endpoints;
\end{definition}
\begin{proof}
  This is the declared model definition of Matches Baseball Throw Scenario.
\end{proof}
\begin{definition}[Has Physical Baseball Throw Parameters]
  \label{def:physics:phyx-mini-0693:phyxminiproblems-problemphyxmini0693-hasphysicalbaseballthrowparameters}
  \lean{PhyXMiniProblems.ProblemPhyXMini0693.HasPhysicalBaseballThrowParameters}
  \uses{def:physics:phyx-mini-0693:phyxminiproblems-problemphyxmini0693-lengthinmeters, def:physics:phyx-mini-0693:phyxminiproblems-problemphyxmini0693-durationinseconds, def:physics:phyx-mini-0693:phyxminiproblems-problemphyxmini0693-speedinmeterspersecond, def:physics:phyx-mini-0693:phyxminiproblems-problemphyxmini0693-accelerationinmeterspersecondsquared, def:physics:phyx-mini-0693:phyxminiproblems-problemphyxmini0693-baseballonebouncesetup}
  Positivity and nondegeneracy conditions for the physical setup
\end{definition}
\begin{proof}
  This is the declared model definition of Has Physical Baseball Throw Parameters.
\end{proof}
\begin{definition}[Satisfies Ideal Level Ground Projectile Laws]
  \label{def:physics:phyx-mini-0693:phyxminiproblems-problemphyxmini0693-satisfiesideallevelgroundprojectilelaws}
  \lean{PhyXMiniProblems.ProblemPhyXMini0693.SatisfiesIdealLevelGroundProjectileLaws}
  \uses{def:physics:phyx-mini-0693:phyxminiproblems-problemphyxmini0693-lengthinmeters, def:physics:phyx-mini-0693:phyxminiproblems-problemphyxmini0693-durationinseconds, def:physics:phyx-mini-0693:phyxminiproblems-problemphyxmini0693-speedinmeterspersecond, def:physics:phyx-mini-0693:phyxminiproblems-problemphyxmini0693-accelerationinmeterspersecondsquared, def:physics:phyx-mini-0693:phyxminiproblems-problemphyxmini0693-baseballonebouncesetup}
  Standard ideal-projectile relations for an arc whose launch and landing heights are equal, with no aerodynamic drag and uniform downward gravity: speed magnitude on returning to the launch height equals launch speed; horizontal range is v² sin (2 α) / g; flight time is 2 v sin α / g. They are stated as general relations for all three arcs and contain neither the requested time ratio nor any answer-choice value
\end{definition}
\begin{proof}
  This is the declared model definition of Satisfies Ideal Level Ground Projectile Laws.
\end{proof}
\begin{definition}[one Bounce Flight Time In Seconds]
  \label{def:physics:phyx-mini-0693:phyxminiproblems-problemphyxmini0693-onebounceflighttimeinseconds}
  \lean{PhyXMiniProblems.ProblemPhyXMini0693.oneBounceFlightTimeInSeconds}
  \uses{def:physics:phyx-mini-0693:phyxminiproblems-problemphyxmini0693-durationinseconds, def:physics:phyx-mini-0693:phyxminiproblems-problemphyxmini0693-baseballonebouncesetup}
  Total time in seconds for the two blue legs of the one-bounce throw
\end{definition}
\begin{proof}
  This is the declared model definition of one Bounce Flight Time In Seconds.
\end{proof}
\begin{definition}[no Bounce Flight Time In Seconds]
  \label{def:physics:phyx-mini-0693:phyxminiproblems-problemphyxmini0693-nobounceflighttimeinseconds}
  \lean{PhyXMiniProblems.ProblemPhyXMini0693.noBounceFlightTimeInSeconds}
  \uses{def:physics:phyx-mini-0693:phyxminiproblems-problemphyxmini0693-durationinseconds, def:physics:phyx-mini-0693:phyxminiproblems-problemphyxmini0693-baseballonebouncesetup}
  Flight time in seconds for the green direct throw
\end{definition}
\begin{proof}
  This is the declared model definition of no Bounce Flight Time In Seconds.
\end{proof}
\begin{definition}[one Bounce To No Bounce Time Ratio]
  \label{def:physics:phyx-mini-0693:phyxminiproblems-problemphyxmini0693-onebouncetonobouncetimeratio}
  \lean{PhyXMiniProblems.ProblemPhyXMini0693.oneBounceToNoBounceTimeRatio}
  \uses{def:physics:phyx-mini-0693:phyxminiproblems-problemphyxmini0693-baseballonebouncesetup, def:physics:phyx-mini-0693:phyxminiproblems-problemphyxmini0693-onebounceflighttimeinseconds, def:physics:phyx-mini-0693:phyxminiproblems-problemphyxmini0693-nobounceflighttimeinseconds}
  The dimensionless ratio requested in the question
\end{definition}
\begin{proof}
  This is the declared model definition of one Bounce To No Bounce Time Ratio.
\end{proof}
\begin{definition}[Answer Choice]
  \label{def:physics:phyx-mini-0693:phyxminiproblems-problemphyxmini0693-answerchoice}
  \lean{PhyXMiniProblems.ProblemPhyXMini0693.AnswerChoice}
  Labels of the four numerical answer choices
\end{definition}
\begin{proof}
  This is the declared model definition of Answer Choice.
\end{proof}
\begin{definition}[displayed Ratio]
  \label{def:physics:phyx-mini-0693:phyxminiproblems-problemphyxmini0693-answerchoice-displayedratio}
  \lean{PhyXMiniProblems.ProblemPhyXMini0693.AnswerChoice.displayedRatio}
  \uses{def:physics:phyx-mini-0693:phyxminiproblems-problemphyxmini0693-answerchoice}
  Dimensionless time ratio printed beside each answer label
\end{definition}
\begin{proof}
  This is the declared model definition of displayed Ratio.
\end{proof}
\begin{definition}[recorded Dataset Answer]
  \label{def:physics:phyx-mini-0693:phyxminiproblems-problemphyxmini0693-recordeddatasetanswer}
  \lean{PhyXMiniProblems.ProblemPhyXMini0693.recordedDatasetAnswer}
  \uses{def:physics:phyx-mini-0693:phyxminiproblems-problemphyxmini0693-answerchoice}
  Answer label recorded by the supplied dataset metadata
\end{definition}
\begin{proof}
  This is the declared model definition of recorded Dataset Answer.
\end{proof}
\begin{definition}[Matches Displayed Three Decimal Ratio]
  \label{def:physics:phyx-mini-0693:phyxminiproblems-problemphyxmini0693-matchesdisplayedthreedecimalratio}
  \lean{PhyXMiniProblems.ProblemPhyXMini0693.MatchesDisplayedThreeDecimalRatio}
  \uses{def:physics:phyx-mini-0693:phyxminiproblems-problemphyxmini0693-baseballonebouncesetup, def:physics:phyx-mini-0693:phyxminiproblems-problemphyxmini0693-onebouncetonobouncetimeratio, def:physics:phyx-mini-0693:phyxminiproblems-problemphyxmini0693-answerchoice}
  Agreement with an answer displayed to three decimal places
\end{definition}
\begin{proof}
  This is the declared model definition of Matches Displayed Three Decimal Ratio.
\end{proof}
\begin{definition}[Is Unique Closest Displayed Ratio]
  \label{def:physics:phyx-mini-0693:phyxminiproblems-problemphyxmini0693-isuniqueclosestdisplayedratio}
  \lean{PhyXMiniProblems.ProblemPhyXMini0693.IsUniqueClosestDisplayedRatio}
  \uses{def:physics:phyx-mini-0693:phyxminiproblems-problemphyxmini0693-baseballonebouncesetup, def:physics:phyx-mini-0693:phyxminiproblems-problemphyxmini0693-onebouncetonobouncetimeratio, def:physics:phyx-mini-0693:phyxminiproblems-problemphyxmini0693-answerchoice}
  The selected displayed value is strictly closer than every alternative
\end{definition}
\begin{proof}
  This is the declared model definition of Is Unique Closest Displayed Ratio.
\end{proof}
\begin{lemma}[bounce Launch Angle Sine exact]
  \label{lem:physics:phyx-mini-0693:phyxminiproblems-problemphyxmini0693-bouncelaunchanglesine-exact}
  \lean{PhyXMiniProblems.ProblemPhyXMini0693.bounceLaunchAngleSine_exact}
  \uses{def:physics:phyx-mini-0693:phyxminiproblems-problemphyxmini0693-baseballonebouncesetup, def:physics:phyx-mini-0693:phyxminiproblems-problemphyxmini0693-matchesbaseballthrowscenario, def:physics:phyx-mini-0693:phyxminiproblems-problemphyxmini0693-hasphysicalbaseballthrowparameters, def:physics:phyx-mini-0693:phyxminiproblems-problemphyxmini0693-satisfiesideallevelgroundprojectilelaws}
  Equal total ranges and the half-speed rebound force the lower blue launch angle to satisfy sin θ = 1 / sqrt 5. This is a derived statement, not a premise of the physical model
\end{lemma}
\begin{proof}
  The conclusion follows from the governing relations and figure readouts named in the statement of bounce Launch Angle Sine exact.
\end{proof}
\begin{lemma}[one Bounce To No Bounce Time Ratio exact]
  \label{lem:physics:phyx-mini-0693:phyxminiproblems-problemphyxmini0693-onebouncetonobouncetimeratio-exact}
  \lean{PhyXMiniProblems.ProblemPhyXMini0693.oneBounceToNoBounceTimeRatio_exact}
  \uses{def:physics:phyx-mini-0693:phyxminiproblems-problemphyxmini0693-baseballonebouncesetup, def:physics:phyx-mini-0693:phyxminiproblems-problemphyxmini0693-matchesbaseballthrowscenario, def:physics:phyx-mini-0693:phyxminiproblems-problemphyxmini0693-hasphysicalbaseballthrowparameters, def:physics:phyx-mini-0693:phyxminiproblems-problemphyxmini0693-satisfiesideallevelgroundprojectilelaws, def:physics:phyx-mini-0693:phyxminiproblems-problemphyxmini0693-onebouncetonobouncetimeratio}
  Adding the two blue flight times and comparing with the green 45° flight gives the exact dimensionless ratio 3 / sqrt 10
\end{lemma}
\begin{proof}
  The conclusion follows from the governing relations and figure readouts named in the statement of one Bounce To No Bounce Time Ratio exact.
\end{proof}
% --- Archon declaration topology end ---
% --- Archon physics formalization source end ---
